%!TeX program = xelatex
\documentclass{SYSUReport}


\headl{实验 4：8 位加减法器}
\headr{计算机组成原理实验}
\lessonTitle{《计算机组成原理实验》}
\reportTitle{（实验 4——8 位加减法器） }
\inst{计算机学院}
\major{保密管理}
\stuname{吕梓文}
\stuid{24353028}
\date{\today}

\AddEnumerateCounter{\chinese}{\chinese}{}
\begin{document}

\cover
\thispagestyle{empty} % 首页不显示页码
\clearpage

\pagenumbering{arabic}
\setcounter{page}{1}

%%可选择这里也放一个标题
%\begin{center}
%    \title{ \Huge \textbf{{标题}}}
%\end{center}


\begin{enumerate}[label={\chinese*、},labelsep=0pt]
    \item \textbf{实验目的}
    \begin{enumerate}[label = \arabic*、]
        \item 学习使用Verilog语法进行模块设计；
        \item 通过设计加减法器，加深理解Vivado中组合电路的设计原理和设计方法；
        \item 学习了解使用补码表示解决加减电路一体化的设计原理。
    \end{enumerate}

    \item \textbf{实验内容}
    \begin{enumerate}[label = \arabic*、]
        \item 实现8位带进位、溢出、判零等功能的加减法运算器；
        \item 能正确仿真并在板子上完成设计实现；
        \item 撰写实验报告。
    \end{enumerate}

    \item \textbf{实验原理}
    
    本实验的核心原理包括：
    \begin{enumerate}[label = \arabic*、]
        \item 补码加减法原理
        
        现代数字系统普遍采用补码表示有符号整数。补码的核心性质是：加法、减法都可以用同一套硬件结构实现。

        对于 8 位补码数$a$和$b$：
        \begin{itemize}
            \item 加法：直接执行$a+b$
            \item 减法：执行$a+(-b)$，其中$-b=\lnot b + 1 $
        \end{itemize}
        因此 ALU 中的加减法可以统一为加法器完成，差别仅在于运算模式控制信号。

        \item 标志位（flag）生成原理
        
        为了模拟 CPU 中 ALU 的状态寄存器，本实验生成以下标志位：
        \begin{itemize}
            \item 溢出标志 OF
            
            \textbf{仅用于有符号运算。}
            \begin{figure}[!htbp]
                \centering
                \includegraphics[width = .7\textwidth]{of.png}
                \caption{加减法的溢出情况}
                \label{fig:of}
            \end{figure}

            如图 1 所示：
            \begin{itemize}
                \item  加法溢出条件：两个操作数符号相同，但结果符号发生改变。
                \item  减法溢出条件：两个操作数符号不同，但结果符号与被减数符号不同。
            \end{itemize}
            原理均来源于补码加减法的符号位逻辑。

            \item 符号标志 SF
            
            在有符号模式中，表示结果的符号，即计算结果的最高位。最高位为 1 表示结果为负数，0 表示结果为正数

            \item 零标志 ZF
            
            若结果为 0，则置 1。

            \item 进位/借位标志 CF
            
            \textbf{仅用于无符号运算。}

            将结果用 9 位位宽的寄存器变量存储。当进行无符号运算发生进位或借位时，第 9 位为 1；

        \end{itemize}

    \end{enumerate}

    \item \textbf{实验器材}
    
    电脑一台，Xilinx Vivado 软件一套，Basys3板一块。

    \item \textbf{实验过程与结果}
    
    \begin{enumerate}[label = \arabic*、]
        \item 设计源代码如下：

        \begin{minted}{verilog}
            `timescale 1ns / 1ps
            module alu(
                input clk,
                input [7:0] a,
                input [7:0] b,
                input btnd,
                output reg [7:0] out = 8'b00000000,
                output reg [3:0] flag = 4'b0000,
                output reg compute_mode = 0
            );

                wire btnd_clean;

                debounce d1(
                    .clk(clk),
                    .btn(btnd),
                    .btn_clean(btnd_clean)
                );

                reg [7:0] result8;      // 八位，用来存储计算结果
                reg [8:0] result9;      // 九位，第九位（最高位）用来判断 CF
                reg [3:0] flag_next;

                always @(posedge clk) begin
                    if(btnd_clean)
                        compute_mode <= ~compute_mode;
                end

                // 加减法计算
                always @(*) begin
                    if(compute_mode == 0) begin
                        result9  = {1'b0, a} + {1'b0, b};   
                        result8  = result9[7:0];            
                    end else begin
                        result9  = {1'b0, a} + {1'b0, ~b + 8'b1};
                        result8  = result9[7:0];
                    end
                end

                // 标志位设置
                always @(*) begin
                    // OF
                    if(compute_mode == 0) begin
                        flag_next[0] = (a[7] == b[7]) && (result8[7] != a[7]);
                    end else begin
                        flag_next[0] = (a[7] != b[7]) && (result8[7] != a[7]);
                    end

                    // SF
                    flag_next[1] = result8[7];

                    // ZF
                    flag_next[2] = (result8 == 8'b0);

                    // CF
                    flag_next[3] = result9[8];
                end

                always @(posedge clk) begin
                    out  <= result8;
                    flag <= flag_next;
                end

            endmodule


            // 消抖模块
            module debounce(
                input clk,
                input btn,
                output reg btn_clean
            );

                reg [19:0] cnt = 0;
                reg stable = 0;

                always @(posedge clk) begin
                    if(btn == stable) begin
                        cnt <= 0;
                    end else begin
                        cnt <= cnt + 1;
                        if(cnt == 100000) begin
                            stable <= btn;
                        end
                    end
                end

                reg last;
                always @(posedge clk) begin
                    last <= stable;
                    btn_clean <= (~last & stable);
                end
            endmodule
        \end{minted}
        代码实现说明：
        \begin{enumerate}[label = （\arabic*）、]
            \item 按键消抖（\verb|debounce|模块）
            
            消抖模块由稳定检测、计数判定和上升沿输出三部分组成
            \begin{itemize}
                \item 输入状态稳定性检测

                当按钮输入保持不变时计数清零，若检测到变化则开始计数，达到阈值判定为有效翻转：
                \begin{minted}{verilog}
                    always @(posedge clk) begin
                        if(btn == stable) begin
                            cnt <= 0;
                        end else begin
                            cnt <= cnt + 1;
                            if(cnt == 100000) begin
                                stable <= btn;
                            end
                        end
                    end
                \end{minted}
                
                \item 上升沿检测
                
                通过记录 \verb|last| 生成脉冲作为有效按键事件：
                \begin{minted}{verilog}
                    always @(posedge clk) begin
                        last <= stable;
                        btn_clean <= (~last & stable);
                    end
                \end{minted}
                这确保每次按下按键只触发一次模式切换，避免抖动导致的连跳。
            \end{itemize}
            
            \item ALU 加减法逻辑
            
            \begin{minted}{verilog}
                always @(*) begin
                    if(compute_mode == 0) begin
                        result9  = {1'b0, a} + {1'b0, b};   
                        result8  = result9[7:0];            
                    end else begin
                        result9  = {1'b0, a} + {1'b0, ~b + 8'b1};
                        result8  = result9[7:0];
                    end
                end
            \end{minted}
            \begin{itemize}[leftmargin=0pt]
                \item 两种运算都构建为 9 位加法器输入；
                \item \verb|~b + 1|即为 b 的补码，实现减法；
                \item 9 位结果\verb|result9[8]|用于 CF 判断；
                \item 输出部分始终取 \verb|result9| 的低 8 位。
            \end{itemize}
            \item 标志位 \verb|flag| 的生成
            
            标志位生成逻辑为组合逻辑，重点区分有符号与无符号两种语义。该部分直接体现了 CPU 状态寄存器标志位的生成机制。
            \begin{itemize}
                \item 溢出标志 OF 
                溢出仅与有符号加减法相关；
                \begin{minted}{verilog}
                    if(compute_mode == 0) begin
                        flag_next[0] = (a[7] == b[7]) && (result8[7] != a[7]);
                    end else begin
                        flag_next[0] = (a[7] != b[7]) && (result8[7] != a[7]);
                    end
                \end{minted}
                这里严格遵循补码溢出判定规则：
                \begin{itemize}
                    \item  加法：同号输入、异号输出为溢出
                    \item  减法：异号输入、结果符号变化为溢出
                \end{itemize}

                \item 符号位 SF
                
                仅适用于有符号运算，直接取结果最高位：
                \begin{minted}{verilog}
                    flag_next[1] = result8[7];
                \end{minted}

                \item 零标志 ZF
                结果为0即置位：
                \begin{minted}{verilog}
                    flag_next[2] = (next_out == 0);
                \end{minted}

                \item 无符号进位/借位 CF
                
                仅在无符号语义下有效，从 \verb|result9[8]|
                \begin{minted}{verilog}
                   flag_next[3] = result9[8];
                \end{minted}

                使用 9 位宽度是为了捕获加法进位或减法借位，这是无符号运算不可或缺的标志位。

            \end{itemize}
            \item 结果寄存
            所有运算与标志位均在时钟上升沿锁存输出，符合同步电路规范：
            \begin{minted}{verilog}
                always @(posedge clk) begin
                    out  <= result8;
                    flag <= flag_next;
                end
            \end{minted}
        \end{enumerate}
    \item 约束源代码如下：
    \begin{minted}{tcl}
        # 时钟绑定
        set_property -dict { PACKAGE_PIN W5 IOSTANDARD LVCMOS33 } [get_ports clk]
        create_clock -period 10.000 -name clk -waveform {0 5} [get_ports clk]

        # 第一个操作数 a
        set_property PACKAGE_PIN V17 [get_ports a[0]]
        set_property IOSTANDARD LVCMOS33 [get_ports a[0]]
        set_property PACKAGE_PIN V16 [get_ports a[1]]
        set_property IOSTANDARD LVCMOS33 [get_ports a[1]]
        set_property PACKAGE_PIN W16 [get_ports a[2]]
        set_property IOSTANDARD LVCMOS33 [get_ports a[2]]
        set_property PACKAGE_PIN W17 [get_ports a[3]]
        set_property IOSTANDARD LVCMOS33 [get_ports a[3]]
        set_property PACKAGE_PIN W15 [get_ports a[4]]
        set_property IOSTANDARD LVCMOS33 [get_ports a[4]]
        set_property PACKAGE_PIN V15 [get_ports a[5]]
        set_property IOSTANDARD LVCMOS33 [get_ports a[5]]
        set_property PACKAGE_PIN W14 [get_ports a[6]]
        set_property IOSTANDARD LVCMOS33 [get_ports a[6]]
        set_property PACKAGE_PIN W13 [get_ports a[7]]
        set_property IOSTANDARD LVCMOS33 [get_ports a[7]]

        # 第二个操作数 b
        set_property PACKAGE_PIN V2 [get_ports b[0]]
        set_property IOSTANDARD LVCMOS33 [get_ports b[0]]
        set_property PACKAGE_PIN T3 [get_ports b[1]]
        set_property IOSTANDARD LVCMOS33 [get_ports b[1]]
        set_property PACKAGE_PIN T2 [get_ports b[2]]
        set_property IOSTANDARD LVCMOS33 [get_ports b[2]]
        set_property PACKAGE_PIN R3 [get_ports b[3]]
        set_property IOSTANDARD LVCMOS33 [get_ports b[3]]
        set_property PACKAGE_PIN W2 [get_ports b[4]]
        set_property IOSTANDARD LVCMOS33 [get_ports b[4]]
        set_property PACKAGE_PIN U1 [get_ports b[5]]
        set_property IOSTANDARD LVCMOS33 [get_ports b[5]]
        set_property PACKAGE_PIN T1 [get_ports b[6]]
        set_property IOSTANDARD LVCMOS33 [get_ports b[6]]
        set_property PACKAGE_PIN R2 [get_ports b[7]]
        set_property IOSTANDARD LVCMOS33 [get_ports b[7]]


        # 运算结果
        set_property PACKAGE_PIN U16 [get_ports out[0]]
        set_property IOSTANDARD LVCMOS33 [get_ports out[0]]
        set_property PACKAGE_PIN E19 [get_ports out[1]]
        set_property IOSTANDARD LVCMOS33 [get_ports out[1]]
        set_property PACKAGE_PIN U19 [get_ports out[2]]
        set_property IOSTANDARD LVCMOS33 [get_ports out[2]]
        set_property PACKAGE_PIN V19 [get_ports out[3]]
        set_property IOSTANDARD LVCMOS33 [get_ports out[3]]
        set_property PACKAGE_PIN W18 [get_ports out[4]]
        set_property IOSTANDARD LVCMOS33 [get_ports out[4]]
        set_property PACKAGE_PIN U15 [get_ports out[5]]
        set_property IOSTANDARD LVCMOS33 [get_ports out[5]]
        set_property PACKAGE_PIN U14 [get_ports out[6]]
        set_property IOSTANDARD LVCMOS33 [get_ports out[6]]
        set_property PACKAGE_PIN V14 [get_ports out[7]]
        set_property IOSTANDARD LVCMOS33 [get_ports out[7]]

        # 标志位
        set_property PACKAGE_PIN V13 [get_ports flag[0]]
        set_property IOSTANDARD LVCMOS33 [get_ports flag[0]]
        set_property PACKAGE_PIN V3 [get_ports flag[1]]
        set_property IOSTANDARD LVCMOS33 [get_ports flag[1]]
        set_property PACKAGE_PIN W3 [get_ports flag[2]]
        set_property IOSTANDARD LVCMOS33 [get_ports flag[2]]
        set_property PACKAGE_PIN U3 [get_ports flag[3]]
        set_property IOSTANDARD LVCMOS33 [get_ports flag[3]]

        # 模式切换
        set_property PACKAGE_PIN P3 [get_ports compute_mode]
        set_property IOSTANDARD LVCMOS33 [get_ports compute_mode]


        # 按键绑定
        set_property -dict { PACKAGE_PIN U17 IOSTANDARD LVCMOS33 } [get_ports btnd]

    \end{minted}
    \item 运行仿真验证
    
    \begin{figure}[!htbp]
        \centering
        \includegraphics[width = .7\textwidth]{sim.jpg}
        \caption{仿真运行结果}
        \label{fig:sim}
    \end{figure}

    \begin{itemize}
        \item $127 + (-1)$，运算结果为\verb|8'b01111110|（$126$），发生了进位，标志位只有 CF 为 1；
        \item $127 - (-1)$，运算结果为\verb|8'b10000000|（$-128$），发生了上溢，标志位 OF 与 SF 为 1；
        \item $-128 + (-128)$，运算结果为\verb|8'b00000000|（$0$），发生了进位，溢出，标志位 CF、OF 和 ZF 都为 1；
        \item $-128 - (-128)$，运算结果为\verb|8'b00000000|（$0$），发生了进位，溢出，标志位 CF、OF 和 ZF 都为 1；
    \end{itemize}
    \item 烧板验证
    
    \begin{figure}[!htbp]
        \centering
        \includegraphics[width = .7\textwidth]{127+-1.jpg}
        \caption{$127 + (-1) = 126$的运算结果}
        \label{fig:127pm1}
    \end{figure}
    \begin{figure}[!htbp]
        \centering
        \includegraphics[width = .7\textwidth]{127--1.jpg}
        \caption{$127 - (-1) = 128$的运算结果（有符号溢出）}
        \label{fig:127mm1}
    \end{figure}

    \begin{figure}[!htbp]
        \centering
        \includegraphics[width = .7\textwidth]{-128+-128.jpg}
        \caption{$-128 + (-128) = 0$的运算结果（有符号溢出）}
        \label{fig:m128pm128}
    \end{figure}
    \begin{figure}[!htbp]
        \centering
        \includegraphics[width = .7\textwidth]{-128--128.jpg}
        \caption{$-128 - (-128) = 0$的运算结果}
        \label{fig:m128mm128}
    \end{figure}
    \end{enumerate}
        
\end{enumerate}

\clearpage
\end{document}
